\documentclass[14pt]{extarticle}
\usepackage{preamble}
\geometry{letterpaper, landscape, margin=0.5in}
\usepackage{qrcode}
\renewcommand{\answer}[1]{}
\setlength{\parskip}{4pt}

\begin{document}
\pagestyle{empty}

\daybanner{Unit 4 \textperiodcentered\ Topic 4.4 \textperiodcentered\ TODAY'S PROBLEM}{Two Measurements per Technician}

\noindent\begin{minipage}[t]{0.57\linewidth}
\vspace{0pt}
Suppose a company takes an SRS of 18 of its 300 technicians. Each completes the same task with old and new interfaces; interface order is randomly assigned. The plot shows $d=\text{new time}-\text{old time}$. Researchers ask whether the new interface reduces the population mean time.\par\smallskip \textbf{Nia:} ``The technician seems to be the independent unit, so I would build one difference per person; the subtraction order matters for the tail.''\par \textbf{Colin:} ``There are 36 recorded times, so I would use two independent samples of 18 and a two-sided alternative.''\par
\end{minipage}\hfill
\begin{minipage}[t]{0.40\linewidth}
\vspace{0pt}
\begin{center}
\begin{tikzpicture}
\begin{axis}[
  width=0.77\linewidth,
  height=1.44in,
  ymin=0, ymax=2, ytick=\empty, axis y line=none, axis x line=bottom,
  xlabel={New-interface time minus old-interface time (seconds)}, tick label style={font=\small}, label style={font=\small}
]
\addplot[only marks, mark=*, mark size=2.2pt, color=dayblue] coordinates {
  (-4.2,1)
  (-3.6,1)
  (-3.2,1)
  (-2.9,1)
  (-2.6,1)
  (-2.4,1)
  (-2.2,1)
  (-2.0,1)
  (-1.9,1)
  (-1.7,1)
  (-1.6,1)
  (-1.4,1)
  (-1.2,1)
  (-1.0,1)
  (-0.7,1)
  (-0.4,1)
  (0.0,1)
  (0.6,1)
};
\end{axis}
\end{tikzpicture}
\end{center}
\end{minipage}

\begin{firsttakebox}{FIRST TAKE (5 min, on your sheet)}
Would you compare each technician's two times or treat them as separate people? Explain in one sentence.
\end{firsttakebox}

\begin{description}[leftmargin=1.15in, labelwidth=1in, itemsep=2pt, topsep=2pt]
  \item[\dokbadge{1}~(a)] Define $\mu_d$ in context.
  \item[\dokbadge{2}~(b)] Check the random-sample, 10 percent, and small-sample shape conditions in three short notes.
  \item[\dokbadge{3}~(c)] In 3--4 sentences, choose a setup and justify it. State the hypotheses using the given subtraction order and reduction question, then explain why the 36 times are not 36 independent people.
\end{description}

\vfill
\noindent\begin{minipage}{0.84\linewidth}
\textbf{Video + follow-along:} \href{https://robjohncolson.github.io/apstats-live-worksheet/u7_lesson4_live.html}{\texttt{\detokenize{u7_lesson4_live.html}}} (scan the code)\par
\textbf{Finish (a)--(c) after the video. Turn in your sheet.}
\end{minipage}\hfill
\qrcode[height=0.75in]{https://robjohncolson.github.io/apstats-live-worksheet/u7_lesson4_live.html}

\end{document}
